\documentclass[11pt,letterpaper]{article}
\usepackage{../styles/mae2250-handout}

\renewcommand{\HandoutNumber}{4}
\renewcommand{\HandoutTitle}{Coffee Grinder}
\renewcommand{\PartTitle}{Part 2: Modelling}
\renewcommand{\DueDate}{\{3 days before the next lab section, or enough time for TAs to order parts\}}
\renewcommand{\TeamType}{Individual}

\begin{document}
\MakeHandoutHeader

One part of this assignment is to design test housing for your coffee grinder. The purpose of test housing is to have an easy way to test only a desired mechanism rather than an entire prototype. You will create something similar for Prototype 1 of your ODP. For this housing, you may not use any ordered parts, and the materials do not have to follow budget or space constraints.

\PartHeading{PART 2}
\begin{AssignmentSteps}
  \item \underline{Find parts.}
  \begin{itemize}
    \item On McMaster, locate and download any parts you want to use.
    \begin{itemize}
      \item Be sure to document any changes you make to your design as you see what is available to you.
    \end{itemize}
    \item As you add parts to your design, begin making a bill of materials. This will be how the staff knows exactly what parts to get you, so make sure it's clear. At minimum, it should have:
    \begin{itemize}
      \item A description of each part with relevant parameters (eg diameter, length)
      \item Some method to find the part (eg McMaster serial number, physical location if in TDS)
      \item The part's price
    \end{itemize}
    \item This step will likely be done in tandem with the next.
  \end{itemize}

  \item \underline{CAD the internal mechanism.}
  \begin{itemize}
    \item The gears should be motion-linked so they rotate together. Use the following video to learn how to do so: \{link to last year's video\}
    \item Highlight the locations of your bearings. You may do this however you want, but you may also use the following method:
    \begin{itemize}
      \item Go to Inspect\(\rightarrow\)Display Component Colors. This should turn all of your components into bright pastel colors.
      \item Right click your bearings and click ``Cycle Component Color.'' Cycle through until all of your bearings share the same unique color.
      \begin{itemize}
        \item Please make sure colors are distinct. For example , if one part is light purple, do not turn your bearings a light pink.
      \end{itemize}
      \item Screenshot this setup and return your layout to its original color.
    \end{itemize}
    \item Remember to use the center of gravity test to see where your supports will need to be.
  \end{itemize}

  \item \underline{CAD test housing.}
  \begin{itemize}
    \item Make a sketch of your test housing.
    \item Make another assembly, and add your assembly of the internal mechanism into this.
    \item Add your housing over the mechanism.
  \end{itemize}
\end{AssignmentSteps}

\newpage

\textbf{Move your model and all parts to a folder called ``coffeegrinder\_Final'' and submit a document with your finalized bill of materials, the highlighted bearing locations (if more than two colors in your design, be sure to specify which colors are the bearings), and any updated sketches on Canvas by \{Part 2 due date\}.}

In your next lab section, you will build the test housing you created. If there are any parts that need to come from the RPL, ensure that you place an order with ample time for the RPL to process!

\newpage

\{assuming the students wouldn't have had time to see initial feedback before submitting this part\}

\PartHeading{TA CHECKLIST}
\begin{Checklist}
  \item All pieces in the bill of materials match the internal CAD.
  \item The bill of materials does not exceed \$100.
  \item The bearing locations are obvious in the screenshot.
  \item A pair of gears and a conical burr was used in the design.
  \item The gears are properly motion-linked.
  \item All parts are properly constrained.
  \item The CAD follows good practices.
\end{Checklist}

\end{document}
